\section{Introduzione}

\begin{itemize}
    \item \textbf{Sistema informativo}: è l’insieme delle attività umane e dei dispositivi di memorizzazione ed elaborazione che organizzano e gestiscono le informazioni di interesse per un’organizzazione, di qualunque dimensione.  
    Un sistema informativo \textbf{non contiene necessariamente tecnologia informatica}.
    
    \item \textbf{Dato}: è un elemento grezzo, un fatto isolato o un valore che, preso da solo, non ha significato.
    
    \item \textbf{Informazione}: nasce quando un dato viene interpretato e collocato in un contesto che lo rende comprensibile e utile.
\end{itemize}

\noindent
Ad esempio: il numero \emph{37} da solo è un semplice dato. Se però diciamo che “37 °C è la temperatura corporea normale di un essere umano”, allora stiamo fornendo un’informazione.

\begin{tcolorbox}[colback=white,colframe=red!80!black,title=Per capire meglio]
Un dato è come un mattone: preso da solo non serve a molto.  
Quando i mattoni vengono messi insieme per costruire una casa, essi diventano informazione.
\end{tcolorbox}

\noindent
Per capire meglio come funziona un sistema informativo (cioè come circolano i dati e chi li utilizza) si utilizzano \textbf{diagrammi di flusso}, analizzando quattro aspetti principali:
\begin{enumerate}
    \item Definizione di archivi e sorgenti di dati
    \item Definizione degli utenti
    \item Definizione di procedure e processi
    \item Definizione dei flussi di dati
\end{enumerate}

\begin{itemize}
    \item \textbf{Base di dati}: è una collezione di dati utilizzata per rappresentare, con l’ausilio della tecnologia informatica, le informazioni di interesse di un sistema informativo.
\end{itemize}

\subsection{Problemi e soluzione}

Negli anni ’70 ogni programma aveva i propri file di dati separati, gestiti direttamente dal \textbf{File System}.  
Questa soluzione presentava alcuni problemi:
\begin{itemize}
    \item Scarsa efficienza nell’accesso ai dati su file
    \item Ridondanza dei dati (duplicazione)
    \item Inconsistenza dei dati (informazioni non aggiornate ovunque)
    \item Progettazione dei dati replicata per ogni programma
\end{itemize}

\noindent
Per risolvere questi limiti vennero introdotti i \textbf{DBMS (Database Management System)}: software che fanno da intermediari tra programmi e dati.

\textbf{Vantaggi dei DBMS:}
\begin{itemize}
    \item Maggiore astrazione e potenza espressiva per descrivere i dati
    \item Operazioni complesse di accesso e interrogazione tramite linguaggi specifici
\end{itemize}

\subsection{DBMS}

\begin{itemize}
    \item \textbf{Linguaggi di interazione}:
    \begin{itemize}
        \item \textbf{DDL (Data Definition Language)}: definisce la struttura del database.
        \item \textbf{DML (Data Manipulation Language)}: consente di interrogare e modificare i dati (es. SQL).
    \end{itemize}
\end{itemize}

\begin{tcolorbox}[colback=yellow!20!white, colframe=black, title=Differenza tra Database e DBMS]
\textbf{Database}: è l’insieme organizzato dei dati.  
\emph{Esempio}: un archivio con le informazioni sugli studenti di una scuola.  

\textbf{DBMS (Database Management System)}: è il software che permette di creare, gestire e consultare un database.  
\emph{Esempi}: MySQL, PostgreSQL, Oracle, SQLite.
\end{tcolorbox}

\subsection{Architettura di un DBMS}

\begin{itemize}
    \item \textbf{Schema logico}: rappresentazione delle strutture e delle proprietà della base di dati, definita tramite i costrutti del modello del DBMS.
    \item \textbf{Schema interno}: rappresentazione della base di dati attraverso le strutture fisiche di memorizzazione.
    \item \textbf{Schema esterno}: descrive una porzione dello schema logico di interesse per uno specifico utente o applicazione.
\end{itemize}

\subsection{Indipendenza dei dati}

\begin{itemize}
    \item \textbf{Indipendenza fisica}: lo schema logico della base di dati è indipendente dallo schema interno.
    \item \textbf{Indipendenza logica}: gli schemi esterni della base di dati sono indipendenti dallo schema logico.
\end{itemize}
